\section{Source Section III.B: Adiabatic Expansion Scenario}
\label{sec:adiabatic-expansion}

\sourceeqrange{Scenario figures 1--2; Eq. (3.12)}

\subsection*{Purpose}

This subsection helps the reader locate which source assumptions govern the
adiabatic-expansion example.  It does not reproduce snapshots, time histories,
or stability thresholds.

\subsection*{Reading Guide}

The setup begins from a spatially uniform supercritical state and changes the
wall temperature.  The important bookkeeping objects are the conserved mean
number density, no-slip walls, the liquid-wetting density boundary condition,
and the common wall-temperature branch.  The qualitative mechanism is read
through the balance laws from Section II.D and the wall data from Section III.A,
while Eq. (3.12) supplies the inverse isothermal-compressibility diagnostic.

\subsection*{Finite Estimate}

Equation (3.12) is a local thermodynamic estimate.  The companion records it as
a scalar field used to read the source plot at the cell center; it is not a
numerical reconstruction.

The diagnostic is the inverse isothermal compressibility:
\begin{equation}
  K_T^{-1}
  =n\left(\frac{\partial p}{\partial n}\right)_T .
  \sourceeqbadge{3.12}
  \label{eq:iii-adiabatic-expansion-compressibility-type}
\end{equation}
Using the Section II.A van der Waals pressure branch and differentiating at
fixed \(T\),
\begin{equation}
  \left(\frac{\partial p}{\partial n}\right)_T
  =\frac{k_BT}{(1-v_0n)^2}-2\epsilon v_0n.
\end{equation}
The spinodal-temperature definition gives
\begin{equation}
  T_s=\frac{2\epsilon v_0n(1-v_0n)^2}{k_B}
  \quad\Longrightarrow\quad
  2\epsilon v_0n=\frac{k_BT_s}{(1-v_0n)^2}.
\end{equation}
Therefore
\begin{align}
  K_T^{-1}
  &=n\left[\frac{k_BT}{(1-v_0n)^2}
      -\frac{k_BT_s}{(1-v_0n)^2}\right] \\
  &=\frac{k_Bn}{(1-v_0n)^2}(T-T_s).
  \sourceeqbadge{3.12}
  \label{eq:iii-inverse-compressibility-derived}
\end{align}
As a type check, \((\partial p/\partial n)_T\) has pressure per number-density
units, so multiplication by \(n\) gives pressure.  The sign of \(K_T^{-1}\)
helps the reader locate stable, metastable, and spinodal portions of the source
time trace.  It does not supply an evolution equation for the domain pattern or
the wall-layer growth.

The calculation chain is local: equation (3.12) is evaluated from the local
thermodynamic branch, then read against the cell-center state reported by the
source.  It is not integrated over the cell and it does not determine the
interface pattern.
The finite type and scale checks for this scenario are covered by the
Section III mirror pair listed in Section~\ref{sec:numerical-results-overview};
they do not reproduce the source time trace.

The branch data are:
\[
\begin{array}{c|c}
\text{input} & \text{role} \\
\hline
\langle n\rangle & \text{conserved mean density in the closed cell} \\
T_{\rm wall}(t) & \text{common wall-temperature change} \\
\text{Eq. (3.2)} & \text{liquid-wetting density boundary condition} \\
K_T^{-1} & \text{cell-center thermodynamic diagnostic}
\end{array}
\]

\subsection*{Limitations}

The paper's figures are simulation outcomes.  This companion does not infer
mesh-update rules, time integration, fluctuation effects, or any source-code
choice beyond the printed equations and assumptions.
